% Two stop categories in time, and the latch that outlasts both.
\begin{tikzpicture}[font=\sffamily\scriptsize, x=1mm, y=1mm]

% ================================================================ category 0
\node[chlabel, anchor=west] at (0,42) {category 0: power removed immediately, an uncontrolled stop};
\node[signame] at (-2,34) {velocity};
\draw[taxis] (0,26) -- (124,26);
\draw[signal, line width=0.9pt] (0,36) -- (28,36) -- (34,26) -- (124,26);
\draw[release, draw=red!70!black] (28,22) -- (28,40);
\node[tlabel, anchor=west, text=red!60!black] at (30,38) {the trigger};

\node[signame] at (-2,16) {power};
\draw[signal, line width=0.9pt] (0,20) -- (28,20) -- (28,14) -- (124,14);
\node[tlabel, anchor=west] at (30,17) {removed at once, and the axis stops however it stops};
\node[tlabel, anchor=west, text width=52mm, align=left] at (0,6)
  {\textbf{A safe state for a horizontal or non-backdrivable axis, and the
   opposite of one for an axis that would descend.}};

% ================================================================ category 1
\node[chlabel, anchor=west] at (0,-4) {category 1: a controlled stop with power available, then power removed};
\node[signame] at (-2,-12) {velocity};
\draw[taxis] (0,-20) -- (124,-20);
\draw[signal, line width=0.9pt] (0,-10) -- (28,-10)
   .. controls (44,-10) and (52,-20) .. (68,-20) -- (124,-20);
\draw[release, draw=red!70!black] (28,-24) -- (28,-6);
\node[tlabel, anchor=west, text=red!60!black] at (30,-7.5) {the same trigger};

\node[signame] at (-2,-30) {power};
\draw[signal, line width=0.9pt] (0,-26) -- (68,-26) -- (68,-32) -- (124,-32);
\draw[tspan] (28,-35) -- node[tlabel]{power still available: this is the deceleration} (68,-35);
\node[tlabel, anchor=west] at (70,-29) {and only now removed, with the brake holding};

% ================================================================ the latch
\node[chlabel, anchor=west] at (0,-42) {and the latch, which follows both and outlasts its cause};
\node[fieldbox, fill=red!25, minimum width=86mm, text width=84mm, minimum height=6mm,
      font=\sffamily\tiny] at (28,-52) {latched, cause recorded, lifetime count incremented};
\node[fieldbox, fill=black!8, minimum width=28mm, text width=26mm, minimum height=6mm,
      font=\sffamily\tiny] at (0,-52) {running};
\node[fieldbox, fill=usrcol, minimum width=10mm, text width=8mm, minimum height=6mm,
      font=\sffamily\tiny] at (114,-52) {check};
\draw[release] (28,-56) -- (28,-44);
\draw[release] (114,-56) -- (114,-44);
\node[tlabel, anchor=north] at (28,-57) {trigger};
\node[tlabel, anchor=north] at (100,-57) {an intentional reset, which starts nothing};
\node[tlabel, anchor=south] at (60,-45.5) {the cause goes away somewhere in here, and nothing changes};

\node[umlnote, text width=56mm, anchor=north west] at (0,-64)
  {The reset requires an intentional human action, runs a full self-check, and
   \textbf{does not itself initiate a restart}. That last requirement is in the
   same published clause as the reset, and it is the one most often left out of
   an implementation.};
\node[umlnote, text width=56mm, anchor=north west] at (66,-64)
  {On this bench the lower row is drawn, tested in logic, and never felt: there
   is no brake and no motor, so the deceleration is a computed trajectory and
   the power trace is a pin. The chapter says so wherever the figure appears.};
\end{tikzpicture}
